\noindent\textbf{Einheit 4 · Grundwert berechnen}

\erg{23}{a) Prozentsatz \quad b) Prozentwert \quad c) Grundwert \quad d) Prozentwert \quad e) Grundwert}
\erg{24}{a) 40 Stück \quad b) $70\,€$ \quad c) $40\,$kg \quad d) $48\,$m}
\erg{25}{a) $28\,€$ \quad b) $36\,$kg \quad c) $35\,$m \quad d) $60\,€$ \quad e) $45\,$kg \quad f) $350\,€$}
\erg{26}{a) $1\,\% = 18 : 60 = 0{,}3$; $100\,\% = 30$ Kinder \quad b) $1\,\% = 27\,€ : 15 = 1{,}80\,€$; $100\,\% = 180\,€$}
\erg{27}{a) Tim hat den Prozentsatz vom Prozentwert genommen, statt auf das Ganze zurückzurechnen; richtig ist $35 : 25 = 1{,}4$ und $100 \cdot 1{,}4 = 140$ Gäste \quad b) 300 Karten}
\erg{28}{a) Der Grundwert sind $100\,\%$; aus einem Prozent bekommt man ihn, indem man es hundertmal nimmt \quad b) ja – der Prozentwert ist dann nur ein Teil des Ganzen, und ein Teil ist kleiner als das Ganze}
\erg{29}{a) Prozentwert, 6 Eier \quad b) Prozentsatz, $25\,\%$ \quad c) Grundwert, $60\,€$ \quad d) Prozentwert, 54 \quad e) Grundwert, 50 Kinder}
